The relation between entanglement and boundary area arose in quantum-field
calculations motivated by black-hole entropy.  Bombelli, Koul, Lee, and
Sorkin related entropy from tracing inaccessible quantum-field degrees of
freedom to black-hole area~\cite{BombelliEtAl1986}; Srednicki subsequently found area scaling
numerically for a spherical region in the vacuum of a regulated free
massless scalar field~\cite{Srednicki1993}.  These calculations motivated
the entanglement area-law question through regulated field models. For
finite-dimensional lattice spins, the gap-to-area conjecture asks whether
a unique ground state has an entropy bound proportional to boundary area
when local dimension, interaction range and strength are bounded and the
spectral gap has a size-independent positive lower
bound~\cite{Hastings2007,AradLandauVazirani2012}.
Theorem~\ref{thm:area} gives an affirmative answer in the stated
two-dimensional induced-domain setting.

In one dimension, Hastings proved a length-independent entropy bound for
the unique ground state of a finite-range Hamiltonian with uniformly
bounded interactions, fixed local dimension, and a positive spectral
gap~\cite{Hastings2007}.  Arad, Landau, and Vazirani improved the dependence
on the inverse gap for frustration-free chains using approximate ground-state
projectors~\cite{AradLandauVazirani2012}.  These are operators that preserve
the ground vector while contracting its orthogonal complement, with
controlled growth of Schmidt rank.  Arad, Kitaev, Landau, and Vazirani
extended this approach to general, possibly frustrated
chains~\cite{AradKitaevLandauVazirani2013}.  A complementary result of Brand\~ao and Horodecki
derives an area law for pure one-dimensional states from exponential
decay of correlations between arbitrary separated
regions~\cite{BrandaoHorodecki2015}.  This region-wise hypothesis is
stronger than decay of single-site correlation functions.  More generally,
the exponential clustering theorem relates a spectral gap to decay of
connected correlations of separated
observables~\cite{NachtergaeleSims2006}; an entropy estimate requires
control of the reduced state's full spectrum.
Anshu, Harrow, and Soleimanifar obtained boundary-dependent bounds on
smooth entanglement spread, which measures the width of the significant
part of that spectrum, including square-root boundary scaling on lattices
up to logarithmic factors~\cite{AnshuHarrowSoleimanifar2022}.  Such a width
estimate leaves the location of the spectrum, and hence its entropy,
to be controlled.

Strict area-law bounds in higher dimensions are also known for
restricted model classes.  Cramer, Eisert, Plenio, and Drei\ss ig proved
such a bound for arbitrary regions in cubic harmonic lattices with a
finite-range positive potential matrix whose condition number is
uniformly bounded~\cite{CramerEisertPlenioDreissig2006}.  These Gaussian
models have infinite-dimensional oscillator sites and provide a
comparison with the finite-dimensional spin setting considered here.

Other higher-dimensional results make additional spectral or thermodynamic
assumptions.  Masanes obtained bounds with logarithmic corrections from correlation
decay and bounds on the number of low-energy states of restricted
Hamiltonians~\cite{Masanes2009}; Brand\~ao and Cramer obtained a logarithmically
corrected bound for cubic regions under a low-temperature heat-capacity
hypothesis~\cite{BrandaoCramer2015}.  Cho proved strict area scaling for
balls along boundary-growth families satisfying uniform gaps and an
additional compatibility condition on successive reduced ground
states~\cite{Cho2014}.  These conditions supply information beyond the gap
of the single full Hamiltonian considered here.

For two-dimensional spin systems, Anshu, Arad, and Gosset proved a
subvolume bound for rectangular regions in a rectangular lattice with
plaquette projector interactions, a unique frustration-free ground
state, and a uniform local gap~\cite{AnshuAradGosset2022Subvolume}.
At fixed local dimension and local gap, the bound is
\(O(b^{5/3}\log^{7/3}(2+b))\), where \(b\) is the rectangle's boundary
size.  Their subsequent result gives
\[
  S\le n^{\,1+O((\log n)^{-1/5})}
\]
across a vertical cut of length \(n\), again for locally gapped,
frustration-free plaquette projector models with a unique ground
state~\cite{AnshuAradGosset2022Area}.  The local-gap assumption concerns
restrictions of the Hamiltonian; it is stronger than a gap for the full
system alone.  The latter entropy bound is near-linear, and its theorem
concerns vertical cuts.  The setting considered here allows frustrated
interactions, assumes a gap only for the full Hamiltonian, and includes
arbitrary regions in arbitrary finite induced square-lattice domains.
Arad, Firanko, and Jain extend approximate-projector methods to
maximally mixed ground states and bound their mutual information, with a
two-dimensional application under frustration-free and local-gap
assumptions~\cite{AradFirankoJain2026}.

The proof below combines locality estimates with ways of expressing entropy
through operator norms. For locality, we use the Lieb--Robinson
framework~\cite{LiebRobinson1972,NachtergaeleOgataSims2006} and Kitaev's
centered spectral filter~\cite[Appendix~D.1.2]{Kitaev2006Anyons}.
Swingle and McGreevy discuss positive constraints obtained by squaring
filtered terms, for which retaining a gap is a separate
issue~\cite[Section~8.4]{SwingleMcGreevy2016}. Here taking absolute
values gives the operator domination that retains the full-system gap.
Later, products of positive contractions derived from these constraints
preserve the ground state and suppress excitations, as projection products
do in detectability-lemma arguments~\cite{AnshuAradVidick2016}. Their localization uses propagation
estimates for unital channels, related to dissipative Lieb--Robinson
bounds~\cite{BarthelKliesch2012}; the coupled random-product estimate
needed here is proved in Section~\ref{sec:amplification}.

For the entropy comparison, bounds on local entanglement production are a
precedent for controlling logarithmic commutators independently of
spectator dimensions~\cite{VanAcoleyenMarienVerstraete2013}.
Section~\ref{sec:initial} proves the quadratic real-part estimate needed
for our marginal tails. Strong subadditivity and conditional-entropy
continuity supply the entropy inequalities used throughout
\cite{LiebRuskai1973,AlickiFannes2004,Winter2016}.
We encode marginal spectra on symmetric tensor powers using
Schur--Weyl duality~\cite{EtingofEtAl2011}, whose connection with spectrum
estimation is developed by Keyl and Werner~\cite{KeylWerner2001}.
The passage back to one-copy vectors uses finite quantum de~Finetti
measures~\cite{ChristandlEtAl2007,LewinNamRougerie2015} and the product
calculus of symmetrized finite observables~\cite{RaggioWerner1989}.
Weighted operator geometric means~\cite{PuszWoronowicz1975,KuboAndo1980}
combine the changing partitions, while the modular comparison starts
from the resolvent-compression identity of Carlen and
Vershynina~\cite[Lemma~2.1]{CarlenVershynina2019}.
The detailed construction uses these ingredients to compare entropy gain
and excitation energy in the same family of states; the next overview
explains how that comparison leads to the area bound.
